\chapter{Conclusão}
Esse trabalho propõs-se a avaliar qualitativamente a qualidade do 
código-fonte gerado por modelos de linguagem de grande escala, investigando 
impactos psicológicos sobre os desenvolvedores á luz do viés de automação e do 
avarento cognitivo, fenômenos consolidados na área da saúde antes da 
popularização da inteligência artificial. Para isso adotou-se uma abordagem 
qualitativa, ancorada em critérios de qualidade de normas reconhecidas como a 
ISO/IEC 25010, comparando código escrito manualmente com aquele gerado por 
inteligência artificial. 

As análises preliminares indicaram que a geração de código por inteligência 
artificial dificilmente resulta em um programa correto na primeira tentativa, exigindo 
iterações sucessivas e, sobretudo, a presença de um operador com conhecimento 
aprofundado em programação. Observou-se que a ferramenta trata questões como 
gestão de memória, casos de borda e erros de lógica apenas quando o 
desenvolvedor percebe sua omissão, o que recoloca sobre o humano a 
responsabilidade pela detecção dos defeitos. Esse resultado parcial reforça a pertinência do referencial adotado: o risco não reside apenas na ferramenta, mas na 
aceitação acrítica de seu resultado por quem a utiliza. 

Espera-se que o estudo contribua para uma utilização mais consciente de 
assistentes baseados em inteligência artificial no desenvolvimento de software, 
evidenciando os limites dessas ferramentas e a importância do julgamento crítico do 
desenvolvedor. Como continuidade, prevê-se a ampliação do conjunto de problemas 
e de modelos avaliados, a aplicação sistemática dos critérios de qualidade e o 
aprofundamento da análise qualitativa dos defeitos identificados, de modo a 
consolidar as observações aqui apresentadas.
